\begin{exercise}{24}
Consider the map $[0]\to[n]$ in $\Delta$ with image $\{0\}$.
Show that this determines a map $0\colon\Delta^0\to\partial\Delta^n$.
Calculate the simplicial homotopy sets
$\pi_i^\Delta(\partial\Delta^n,0)$ for $i\ge1$ and $n\ge2$.
Deduce that $\partial\Delta^n$ is not a Kan complex.

\begin{proof}
In the stated range $n\ge2$, the map $[0]\to[n]$ is nonsurjective,
so it is a vertex of $\partial\Delta^n$ by Definition~1.1.37.
Corollary~1.1.12 gives the required map $0$.

A map of pairs $(\Delta^i,\partial\Delta^i)\to(\partial\Delta^n,0)$
is represented by a nonsurjective order-preserving map
$\alpha\colon[i]\to[n]$. Since $i\ge1$, every vertex of
$\Delta^i$ lies in $\partial\Delta^i$, so the boundary condition
forces $\alpha(t)=0$ for every $t\in[i]$.
Thus the constant map is the only map of pairs, and consequently
\[
 \pi_i^\Delta(\partial\Delta^n,0)=\{*\}
 \qquad(i\ge1,\ n\ge2).
\]

If $\partial\Delta^n$ were Kan, Fact~1.1.51 would identify these
sets with the ordinary homotopy groups of its realization.
But $|\partial\Delta^n|\cong S^{n-1}$, as used in the proof of
Corollary~1.1.41, and $\pi_{n-1}(S^{n-1})$ is nontrivial:
the identity of the sphere has degree $1$, whereas nullhomotopic
maps have degree $0$. This contradicts the computation at
$i=n-1\ge1$. Hence $\partial\Delta^n$ is not Kan.
\end{proof}
\end{exercise}
